\section{Experiments}\label{sec:experiments}

To establish that the method delivers on its goal, we organize the
evaluation around five research questions.
\textbf{RQ1} asks whether the full method outperforms state-of-the-art
distillation baselines on standard agent benchmarks under matched
budgets. \textbf{RQ2} asks how much the acquisition component
contributes when the distillation component is held fixed.
\textbf{RQ3} asks how much the distillation component contributes when
every method receives the identical acquired data. \textbf{RQ4} asks
whether each component of the method adds measurable value on its own.
\textbf{RQ5} asks how efficiently the budget is converted into in-task
capability and how off-task behavior follows from that conversion. The
following subsections answer these questions in turn.

\subsection{RQ1, comparison with distillation baselines}
\label{sec:rq1}
Table~\ref{tab:main} reports the main comparison. Every method spends
the same token budget from the same teacher, and the row of deltas
against the untrained base student makes the central pattern visible.
\todo{fill as the BFCL, $\tau^2$-bench, and ALFWorld columns land; on
AppWorld the completed panel already shows every cloning-style
baseline below the base student (0.200--0.215 against 0.248) while
ours is the only method above it (0.284$\pm$.037).} Cloning-style
distillation transfers the teacher's surface style at the expense of
the base model's own competence, and this cost is invisible to
methods that only decide which demonstrations to buy. Our method is
the only one whose training signal preserves the student's verified
behavior while adding the teacher's, and it is correspondingly the
only method that improves on the base student across the panel.

\subsection{RQ2, value of the acquisition component}
\label{sec:rq2}
Table~\ref{tab:cross} fixes our distillation and varies only how the
training data is acquired. \todo{fill from the completed selection
matrix; the controlled-suite column is largely in hand.} Two readings
matter. First, every acquisition baseline improves when paired with
our distillation, confirming that the training-side contribution is
not specific to our own purchases. Second, replacing the baseline
acquisition with our marginal-rule purchasing further improves the
result while cutting the wasted fraction of the budget roughly in
half, because purchases stop at the capability quotas the support set
demands and never flow to out-of-task candidates.

\subsection{RQ3, value of the distillation component}
\label{sec:rq3}
Table~\ref{tab:best} feeds the identical acquired dataset to every
distillation baseline from Table~\ref{tab:main} and to ours. Any
remaining difference is attributable to the training procedure alone.
\todo{fill; existing AppWorld evidence pre-populates most rows: plain
cloning of the same corpus scores 0.197--0.215, an isolated preference
stage collapses, token-selective masking is ineffective, and the
anchored stratified objective reaches 0.284$\pm$.037.} The ordering
matches the mechanism the method is built on. Supervision helps
exactly where the student's own policy fails, and it must arrive
anchored to the student's verified behavior; procedures that clone
indiscriminately, or that push preference pressure without an anchor,
lose the base model's competence before the teacher's is absorbed.

\subsection{RQ4, component ablation}
\label{sec:rq4}
Table~\ref{tab:ablation} builds the method one component at a time,
from plain cloning through the self-anchor with its provenance rule,
the success-axis gated preference term, checkpoint averaging, and
finally marginal acquisition with the displacement cost.
\todo{assemble from the completed AppWorld ladder and the controlled
suite; each step is already measured individually.} Each addition
improves the mean or tightens the variance, and no step requires a
tuned constant, so the full method is the accumulation of individually
verified decisions rather than a jointly tuned pipeline.

\subsection{RQ5, budget conversion and off-task behavior}
\label{sec:rq5}
The final analysis measures where the budget goes.
\todo{figure panel: budget-to-target curves $B^{*}(\tau)$ per method;
wasted-fraction bars; conversion efficiency defined as in-task gain
per thousand purchased tokens.} Our method reaches target performance
at a fraction of the baselines' budgets and wastes 2--4\% of purchased
tokens against roughly half of the strongest generator baseline. The
same accounting explains off-task behavior. On task pairs whose
capability atoms barely overlap, purchases never flow across the
boundary and off-task residue stays at zero at no in-task cost; on
pairs that share response-view atoms, training one side necessarily
moves the other, in the direction the two-view coupling forecast
predicts before any money is spent. A case study
\todo{cot neutralization figure: destructive coupling turned from a
loss increase of +.38 into a decrease of −.40 by ten anchor
demonstrations} shows that a handful of anchor rows neutralizes the
destructive case, closing the loop between the coupling geometry and
the training rule.
